\documentclass[conference]{IEEEtran}
\IEEEoverridecommandlockouts

\usepackage{cite}
\usepackage{amsmath,amssymb,amsfonts}
\usepackage{algorithmic}
\usepackage{graphicx}
\usepackage{textcomp}
\usepackage{xcolor}
\usepackage{booktabs}
\usepackage{array}
\usepackage{url}
\usepackage{hyperref}

\def\BibTeX{{\rm B\kern-.05em{\sc i\kern-.025em b}\kern-.08em
    T\kern-.1667em\lower.7ex\hbox{E}\kern-.125emX}}

\begin{document}

\title{Proposal for CrisisEcho: LLM-Augmented Early Detection\\
of Local Emergencies from Social Media}

\author{
\IEEEauthorblockN{Samuel Enam Zih (Leader)}
\IEEEauthorblockA{
CS, Virginia Tech\\
\texttt{samuelez@vt.edu}}
\and
\IEEEauthorblockN{Christopher Williams}
\IEEEauthorblockA{
CS, Virginia Tech\\
\texttt{cdw24@vt.edu}}
\and
\IEEEauthorblockN{Arjun Sivakumar}
\IEEEauthorblockA{
CS, Virginia Tech\\
\texttt{arjuns8@vt.edu}}
\and
\IEEEauthorblockN{Ki Hong Park}
\IEEEauthorblockA{
CS, Virginia Tech\\
\texttt{stevep00@vt.edu}}
\and
\IEEEauthorblockN{Shiyi Liu}
\IEEEauthorblockA{
CS, Virginia Tech\\
\texttt{lshiyi25@vt.edu}}
}

\maketitle

\begin{abstract}
CrisisEcho is a consumer-facing platform that applies Retrieval-Augmented
Generation (RAG) and multi-step LLM agents to detect, classify, and summarize
hyperlocal emergencies from social media in near real-time. Official emergency
channels routinely lag events by minutes or hours; CrisisEcho closes this gap
by ingesting candidate social and official feeds, storing data across three
MongoDB Atlas databases (document, geospatial, vector) plus Redis and AWS S3,
and running a three-step LangChain agent chain that clusters posts, scores
severity, and generates plain-language alerts. Alerts are delivered via a
Next.js frontend with a Leaflet.js interactive map and WebSocket notifications,
backed by a Go Fiber API. Exact data sources, backend service count, and
non-core libraries are not fixed and will be finalized based on API
accessibility and deployment resources available during development.
\end{abstract}

\begin{IEEEkeywords}
emergency detection, social media mining, RAG, LLM agents, MongoDB Atlas
Vector Search, real-time systems, crisis informatics
\end{IEEEkeywords}

% ─────────────────────────────────────────────────────────────────────────────
\section{Introduction}
% ─────────────────────────────────────────────────────────────────────────────

Official channels --- 911 dispatch, government sensors, police scanners ---
frequently lag fast-moving crises by minutes or hours. Ordinary people post
hyperlocal signals on social media well before any official report: ``Water
rising fast on my street!'' or ``Heard gunshots near the mall.'' Existing
tools are keyword-based and fail on indirect language, sarcasm, and regional
idioms. CrisisEcho replaces rigid keyword matching with semantic RAG-powered
reasoning, demonstrating four LLM-driven paradigm shifts: text to semantics,
retrieval to reasoning, vertical to multi-domain, and closed-world to
open-world generalization.

% ─────────────────────────────────────────────────────────────────────────────
\section{Related Work}
% ─────────────────────────────────────────────────────────────────────────────

Imran et al.~\cite{imran2016twitter} and Olteanu et al.~\cite{olteanu2014crisislex}
established Twitter as a crisis sensor and built foundational lexicons, while
Middleton et al.~\cite{middleton2014real} showed geotagged posts correlate with
ground-truth disaster events. Lewis et al.~\cite{lewis2020rag} introduced RAG
to ground generative models in retrieved evidence, reducing hallucination. Yao
et al.~\cite{yao2023react} demonstrated chained LLM reasoning steps outperform
single-prompt approaches, and Karpukhin et al.~\cite{karpukhin2020dense} showed
dense vector retrieval outperforms BM25 --- motivating our Atlas Vector Search
hybrid retrieval design.

% ─────────────────────────────────────────────────────────────────────────────
\section{Data Sources}
% ─────────────────────────────────────────────────────────────────────────────

\textbf{Note:} The sources below are candidates. The final integrated set will
be determined during development based on API rate limits, access restrictions,
and available deployment resources.

\textbf{Social (candidates):} Twitter/X, Reddit, Bluesky, Mastodon, Nextdoor,
Telegram, Facebook Groups.

\textbf{Official (candidates):} NWS, USGS Earthquake, FEMA, PulsePoint, Waze
Connected Citizens.

\textbf{News/Crowdsourced (candidates):} GDELT, Patch.com RSS, NewsAPI,
Ushahidi, Citizen App.

\textbf{Training/Evaluation:} HumAID~\cite{alam2021humaid} (77k tweets),
CrisisNLP~\cite{imran2016twitter}, CrisisLex~\cite{olteanu2014crisislex},
Kaggle Disaster Tweets, CrisisMMD, FireFlood Dataset.

% ─────────────────────────────────────────────────────────────────────────────
\section{System Architecture}
% ─────────────────────────────────────────────────────────────────────────────

CrisisEcho runs as a five-layer pipeline (Fig.~\ref{fig:architecture}).

\textbf{Layer 1} platform-specific Python workers publish raw JSON to Aiven for
Apache Kafka across four topics (\texttt{social\_raw}, \texttt{official\_alerts},
\texttt{news\_feed}, \texttt{citizen\_reports}).

\textbf{Layer 2} a Python preprocessor applies spaCy cleaning, a cascading
geocoder (Carmen $\rightarrow$ spaCy NER $\rightarrow$ location DB priors),
MinHash LSH deduplication, a DistilBERT relevance filter (removes $\approx$70\%
of posts), Voyage AI text embedding (1024-dim), CLIP image embedding (512-dim),
and S3 image upload.

\textbf{Layer 3} stores data across five backends (Section~\ref{sec:database}).

\textbf{Layer 4} hybrid retrieval feeds a three-step LangChain LLM agent.

\textbf{Layer 5} a Go Fiber API and Next.js frontend deliver alerts.

\begin{figure*}[h]
\centering
\includegraphics[width=\textwidth]{images/CrisisEcho Diagram.pdf}
\caption{CrisisEcho System Architecture}
\label{fig:architecture}
\end{figure*}

% ─────────────────────────────────────────────────────────────────────────────
\section{Database Architecture}
\label{sec:database}
% ─────────────────────────────────────────────────────────────────────────────

MongoDB Atlas M0 (free tier) serves as the unified data platform because a
single Atlas cluster natively supports document storage, 2dsphere geospatial
indexing, and Atlas Vector Search --- no separate vector or geospatial service
required. The free tier is generous for a course project, and Atlas Vector
Search supports multi-modal vector dimensions per collection.

\textbf{DB 1 --- Main (documents + geospatial).} Collections: \texttt{posts}
(unified), per-source raw collections (\texttt{twitter\_posts},
\texttt{reddit\_posts}, etc.), \texttt{clusters}, \texttt{crises},
\texttt{alerts}, \texttt{subscriptions}, \texttt{official\_alerts}. All
location fields carry a 2dsphere index. Key \texttt{posts} fields:
\texttt{post\_id}, \texttt{source}, \texttt{raw\_text}, \texttt{cleaned\_text},
\texttt{timestamp}, \texttt{location} (GeoJSON Point),
\texttt{location\_confidence}, \texttt{text\_embedding\_id},
\texttt{image\_urls} (S3 links), \texttt{cluster\_id}, \texttt{crisis\_type},
\texttt{severity\_score}, \texttt{metadata.username\_hash} (SHA-256).

\textbf{DB 2 --- Location (geocoding support).} Collections:
\texttt{geo\_priors}, \texttt{place\_index}, \texttt{location\_cache}
(TTL-indexed, keyed by SHA-256 text hash). Enables the preprocessing cascade
to avoid re-geocoding identical text.

\textbf{DB 3 --- Vector (Atlas Vector Search).} Collections:
\texttt{text\_embeddings} (1024-dim cosine) and \texttt{image\_embeddings}
(512-dim cosine). Each document carries \texttt{post\_id}, \texttt{timestamp},
\texttt{lat}, \texttt{lng}, \texttt{crisis\_type}, \texttt{is\_relevant} as
filterable metadata inside \texttt{\$vectorSearch}.

\textbf{DB 4 --- AWS S3 (media).} Raw images are uploaded to S3
(\texttt{bucket/<source>/<post\_id>/}); only S3 URLs are stored in MongoDB,
keeping documents lean.

\textbf{DB 5 --- Redis 7 (cache + Pub/Sub).} 30-minute sliding post window,
5-minute LLM summary TTL, and \texttt{alerts:live} Pub/Sub channel relaying
alerts to the Go WebSocket server.

% ─────────────────────────────────────────────────────────────────────────────
\section{RAG Pipeline and Evaluation}
% ─────────────────────────────────────────────────────────────────────────────

\subsection{Hybrid Retrieval}

Each 60-second trigger issues three concurrent queries: (1) Atlas Vector Search
on \texttt{text\_embeddings} (top 50 by cosine similarity, filtered by time
window and bounding box); (2) MongoDB \texttt{\$near} on \texttt{posts}
(geographic proximity regardless of semantics); (3) \texttt{official\_alerts}
lookup injected as high-authority context. Unresolved posts are enriched via
location DB geo priors. Results are merged, deduplicated, and ranked:
$0.5 \times \text{vec} + 0.3 \times \text{recency} + 0.2 \times \text{authority}$.

\subsection{Three-Step LLM Agent}

Claude Haiku (primary) or Ollama Llama3 (fallback) executes: \textbf{Step 1}
cluster posts into distinct events (JSON output); \textbf{Step 2} score
severity 1--5 with rationale; \textbf{Step 3} generate a 2--3 sentence public
alert with no PII. Publish threshold: severity $\geq$ 3, $\geq$ 3 independent
contributors from $\geq$ 2 platforms, geocoding confidence above minimum.

\subsection{Evaluation}

Compared against a CrisisLex keyword baseline on the same corpus. Targets:
Precision@K $>$ 0.75, Recall $>$ 0.70, end-to-end latency $<$ 90s, FPR $<$ 0.20,
geocoding accuracy $>$ 0.80, F1 improvement $>$ 15 pp. Ground truth:
HumAID/CrisisNLP test splits (offline) + PulsePoint/NWS active alerts (online).

% ─────────────────────────────────────────────────────────────────────────────
\section{Technology Stack}
% ─────────────────────────────────────────────────────────────────────────────

Table~\ref{tab:techstack} lists the confirmed core stack. Non-core libraries
and auxiliary tooling are subject to change during development. Backend service
count may be adjusted based on implementation decisions.

\begin{table}[h]
\caption{CrisisEcho Core Technology Stack}
\label{tab:techstack}
\renewcommand{\arraystretch}{1.2}
\begin{tabular}{p{1.4cm} p{2.0cm} p{4.2cm}}
\toprule
\textbf{Component} & \textbf{Technology} & \textbf{Rationale} \\
\midrule
Backend & Go (Fiber) & High throughput; goroutine WebSocket delivery \\
AI Pipeline & Python + LangChain & LLM/NLP ecosystem; LCEL agent chains \\
Frontend & Next.js + Leaflet.js & SSR support; geospatial map \\
Main DB & MongoDB Atlas M0 & 2dsphere geospatial; free cloud tier \\
Location DB & MongoDB Atlas M0 & Geocoding cache + geo priors \\
Vector DB & Atlas Vector Search & Multi-modal ANN; no extra service \\
Media & AWS S3 & Image storage; URLs only in MongoDB \\
Cache & Redis 7 & Post window; LLM TTL; WS relay \\
Queue & Aiven Kafka & Managed Kafka; decouples ingestion \\
LLM & Claude Haiku / Llama3 & Cost-efficient; free offline fallback \\
Deploy & GCP Cloud Run & Serverless; free-tier prototype \\
\bottomrule
\end{tabular}
\end{table}

% ─────────────────────────────────────────────────────────────────────────────
\begin{thebibliography}{00}

\bibitem{imran2016twitter}
M. Imran, P. Mitra, and C. Castillo, ``Twitter as a lifeline: Human-annotated
Twitter corpora for NLP of crisis-related messages,'' in \textit{Proc. LREC},
2016.

\bibitem{olteanu2014crisislex}
A. Olteanu, C. Castillo, F. Diaz, and S. Vieweg, ``CrisisLex: A lexicon for
collecting and filtering microblogged communications in crises,'' in
\textit{Proc. ICWSM}, 2014.

\bibitem{alam2021humaid}
F. Alam, F. Sajjad, H. Dalvi, N. Durrani, K. Darwish, and P. Nakov, ``HumAID:
Human-annotated disaster incidents data from Twitter with deep learning
benchmarks,'' in \textit{Proc. ICWSM}, 2021.

\bibitem{middleton2014real}
S. E. Middleton, L. Middleton, and S. Modafferi, ``Real-time crisis mapping of
natural disasters using social media,'' \textit{IEEE Intelligent Systems},
vol. 29, no. 2, pp. 9--17, 2014.

\bibitem{lewis2020rag}
P. Lewis et al., ``Retrieval-augmented generation for knowledge-intensive NLP
tasks,'' in \textit{Proc. NeurIPS}, 2020.

\bibitem{yao2023react}
S. Yao et al., ``ReAct: Synergizing reasoning and acting in language models,''
in \textit{Proc. ICLR}, 2023.

\bibitem{karpukhin2020dense}
V. Karpukhin et al., ``Dense passage retrieval for open-domain question
answering,'' in \textit{Proc. EMNLP}, 2020.

\end{thebibliography}

% ─────────────────────────────────────────────────────────────────────────────
\appendix
% ─────────────────────────────────────────────────────────────────────────────

\section{Team Role Summary}

\begin{table}[h]
\caption{Team Role Summary}
\label{tab:roles}
\renewcommand{\arraystretch}{1.2}
\begin{tabular}{p{2.0cm} p{5.8cm}}
\toprule
\textbf{Member} & \textbf{Responsibilities} \\
\midrule
S. Zih (Leader) & Backend: AI/DB/API/Docker; Frontend: map; Report: contribution \\
A. Sivakumar & Backend: ingestion/Kafka; Frontend: map UI --- markers, pin drawer; Report: contribution \\
Ki Hong Park & Backend: WebSocket, NLP, Redis; Frontend: query bar, notifications; Report: contribution \\
Shiyi Liu & Backend: evaluation, ablations; Frontend: results display, data QC; Report: proof and submission \\
C. Williams & Frontend: UI/UX, dashboard, diagrams; Backend: deployment, ethics; Report: lead and integration \\
\bottomrule
\end{tabular}
\end{table}

\section{Project Schedule}

\begin{table}[h]
\caption{Project Schedule (13 Weeks)}
\label{tab:schedule}
\renewcommand{\arraystretch}{1.3}
\begin{tabular}{p{1.0cm} p{6.4cm}}
\toprule
\textbf{Weeks} & \textbf{Milestone} \\
\midrule
1--2 & Literature review; provision Atlas, Aiven Kafka, Redis; define schemas \\
3--4 & Ingestion workers; spaCy + geocoding; DistilBERT fine-tune; embeddings to Atlas Vector Search \\
5--6 & Hybrid retrieval; LangChain 3-step agent; Haiku + Llama3 fallback \\
7--8 & Cluster persistence; Redis Pub/Sub; Go Fiber REST + WebSocket API \\
9--10 & Next.js frontend; evaluation suite; keyword baseline comparison \\
11--12 & Stretch features; system testing; latency and threshold tuning \\
13 & Prompt refinement; documentation; demo; final report and submission \\
\bottomrule
\end{tabular}
\end{table}

\section{Acknowledgment}

All architectural and design decisions were made by the project team.

\end{document}
